\documentclass[french]{beamer}
\setbeamertemplate{navigation symbols}{}

\usepackage[utf8]{inputenc}
\usepackage[T1]{fontenc}
\usepackage{graphicx}
\usepackage{tabularx}

\usepackage{xcolor}

\usetheme{Hannover}
\useoutertheme{infolines}
%%%%%%%%
\makeatletter
\beamer@nav@subsectionstyle{show/show/hide}
\makeatother


% \usepackage{siunitx}
% \usepackage[]{algorithm2e}

\usepackage{caption}
\captionsetup{font=scriptsize,labelfont=scriptsize}
\setbeamertemplate{caption}[numbered]

\usepackage{amsmath}
\DeclareMathOperator{\sinc}{sinc}
\DeclareMathOperator{\e}{e}
\setbeamerfont{caption}{size=\footnotesize}


\title[Soutenance MT3]{Etude de la nutation de la terre par analyse de données VLBI}
\author{
Hottier.C, Nurulhuda.I, Robert.C}
\date{4 mars 2016}


\newcommand{\frtt}[1]{\frametitle{#1}}
\newcommand{\FIXME}[1]{\textcolor{red}{FIXME : #1}}

\begin{document}

\begin{frame}
  %\titlepage % Print the title page as the first slide
  \maketitle
  \begin{center}
    \includegraphics[height=0.20\textheight]{Figures/siteon0.jpg}
    %\hspace{8.5cm}
  \end{center}
\end{frame}

\begin{frame}
  \frametitle{Sommaire}
  \tableofcontents
\end{frame}

\section{Introduction}

\begin{frame}
  \frtt{Introduction}

  \begin{itemize}
    \item \'Etude de la nutation de la terre 
    \item Modèle de terre :
      \begin{itemize}
        \item Surface déformable
        \item Manteaux rigide
        \item Noyaux fluide
      \end{itemize}
  \end{itemize}

\end{frame}


\section{Mouvement du pôle et fonction de transfert}

\begin{frame}
  \frtt{Mouvement du pôle et fonction de transfert}
   Ren et Nutation observé relié par une fonction de transfert :
   $$ \tilde{T}(\sigma)  $$
\end{frame}


\section{Données observationnelles}

\begin{frame}
  \frtt{Observation VLBI}
\end{frame}

\section{Analyse des données}

\begin{frame}
  \frtt{Analyse des données}
\end{frame}

\section{Mouvement Libre}
\begin{frame}
  \frtt{Mouvement Libre}
\end{frame}


\section{Conclusion}

\begin{frame}
  \frtt{Conclusion}
\end{frame}

\end{document}
